\documentclass[letterpaper]{article}

\usepackage[margin=1in]{geometry}
\usepackage{amsmath,amssymb,amsfonts}
\usepackage{booktabs}
\usepackage{multirow}
\usepackage{graphicx}
\usepackage{hyperref}
\usepackage{xcolor}
\usepackage{enumitem}
\usepackage{natbib}
\usepackage{authblk}

\title{When Does Multivariate Conformal Prediction Help Robot Action Correction?\\A Systematic Empirical Study}

\author[1]{Anonymous Authors}
\affil[1]{Anonymous Institution}
\date{}

\begin{document}
\maketitle

%==============================================================================
\begin{abstract}
%==============================================================================
Recent work has shown that conformal prediction (CP) can provide calibrated uncertainty estimates for robot action generation.
A natural hypothesis is that \emph{multivariate} CP---which captures per-degree-of-freedom (DOF) uncertainty and cross-dimensional covariance structure---should enable richer, dimension-selective action correction compared to scalar uncertainty scores.
We test this hypothesis systematically by building ACR-VLA, a pipeline that extracts 45-dimensional structured conformal features from flow-matching VLA samples via a Conformal Geometry Score (CGS) and feeds them into a learned Conformal Residual Module (CRM) for inference-time action correction.
Through 11 controlled experiments on LIBERO and LIBERO-Plus benchmarks, we find that multivariate conformal features provide \textbf{no statistically significant benefit} over scalar or zero-dimensional baselines in the practical regime of $K{=}10$ flow-matching samples in 7D action spaces.
The root cause is covariance degeneracy: Ledoit-Wolf shrinkage intensity reaches 0.91 (effective rank 1.3), collapsing multivariate structure to near-diagonal form.
We further show that the self-referential Chunk Overlap Consistency (COC) signal fails as a calibration proxy ($\rho{=}0.18$ with task success), and that inference-time correction provides diminishing returns when stacked on strong training-time methods ($+1.1\%$ on SRPO, $p{=}0.29$).
Our negative results provide precise characterization of when multivariate CP degrades to marginal equivalence, offering practitioner guidance for uncertainty-driven robot action correction.
\end{abstract}

%==============================================================================
\section{Introduction}
\label{sec:intro}
%==============================================================================

Vision-Language-Action (VLA) models based on flow-matching generation have achieved near-saturated performance on standard manipulation benchmarks, yet exhibit catastrophic robustness gaps under distribution shift.
LIBERO-Plus demonstrates this clearly: state-of-the-art VLAs drop from $>$95\% to $\sim$28\% success under realistic perturbations spanning camera, lighting, background, layout, and language variations.

Conformal prediction (CP) offers a principled framework for constructing prediction sets with finite-sample coverage guarantees, requiring only exchangeability rather than distributional assumptions.
In robotic action generation, CP has been applied for anomaly detection, sample selection, and planning-loop integration.
A natural extension is \emph{multivariate} CP: rather than collapsing uncertainty to a scalar, one can construct joint confidence ellipsoids via Mahalanobis-based nonconformity scores, capturing per-DOF uncertainty magnitudes and their covariance structure.

The appeal is clear: if a VLA is highly uncertain about wrist rotation but confident about end-effector position, a correction module conditioned on this structured information could selectively intervene on the uncertain dimensions while preserving confident predictions.
Prior work in low-dimensional settings (3D Minimum Viable Product environments) has validated this intuition, showing $\sim$17\% relative improvement from per-DOF over scalar uncertainty features.

\textbf{This paper asks:} does this advantage transfer to realistic VLA action spaces?
We construct ACR-VLA---Adaptive Conformal Refinement for VLAs---a full pipeline comprising:
(i) a Conformal Geometry Score (CGS) that extracts 45D per-DOF features from $K{=}10$ flow-matching samples via Ledoit-Wolf covariance estimation,
(ii) Adaptive Conformal Inference (ACI) with a self-referential Chunk Overlap Consistency (COC) signal for online threshold adaptation, and
(iii) a learned Conformal Residual Module (CRM) that maps observations and conformal features to action corrections.

Through 11 experiments---a full pipeline evaluation, 7 ablations (A1--A7), per-perturbation analysis, feature importance analysis, and a SnapFlow integration study---we systematically test three claims:
\begin{enumerate}[nosep]
\item Per-DOF conformal features enable dimension-selective correction superior to scalar or no-uncertainty baselines.
\item COC provides a reliable self-referential calibration signal for online adaptation without external supervision.
\item Inference-time conformal correction provides orthogonal value when stacked on training-time methods.
\end{enumerate}

All three claims fail their pre-registered kill criteria.
Rather than abandoning the investigation, we analyze \emph{why} each fails and extract practitioner-relevant insights about the conditions under which multivariate CP can and cannot help.

\paragraph{Contributions.}
(1) We provide a systematic negative result showing that per-DOF conformal features from $K{=}10$ flow-matching samples in 7D action spaces carry negligible information beyond observation embeddings, with precise covariance diagnostics explaining the failure (shrinkage intensity 0.91, effective rank 1.3).
(2) We characterize COC signal failure modes, showing it measures prediction consistency rather than correctness, and quantify the gap between self-referential and supervised calibration signals.
(3) We provide per-perturbation analysis showing that inference-time correction benefits are confined to perturbation types coverable by synthetic augmentation, with negative transfer on unseen types.

%==============================================================================
\section{Background and Related Work}
\label{sec:background}
%==============================================================================

\paragraph{Conformal Prediction for Robotics.}
CP provides distribution-free prediction sets with guaranteed marginal coverage. In robotics, it has been applied for safe planning, anomaly detection, and uncertainty quantification. ReconVLA uses Euclidean distance as a scalar nonconformity score for sample selection among flow-matching outputs. CP-in-the-Loop integrates conformal sets into planning loops. UPS uses CP for policy switching. None of these methods feed structured per-DOF conformal features directly into a learned correction module.

\paragraph{Action Correction at Inference Time.}
A2C2 and PLD learn residual policies conditioned on observations but do not use calibrated uncertainty features. SAFE detects out-of-distribution states for intervention. RobustVLA, DiG-Flow, and SRPO address robustness through training-time data augmentation and RL fine-tuning respectively.

\paragraph{Multivariate CP and Covariance Estimation.}
Extending CP to multivariate outputs requires estimating a joint nonconformity score, typically via Mahalanobis distance with an estimated covariance matrix. With $K$ samples in $p$ dimensions, the sample covariance is rank-deficient when $K < p$ and poorly conditioned when $K/p$ is small. Ledoit-Wolf shrinkage provides a regularized estimator by linearly interpolating between sample covariance and a structured target (scaled identity), with shrinkage intensity $\lambda \in [0,1]$ determined analytically. When $\lambda \to 1$, the estimator collapses to scaled identity, eliminating cross-dimensional structure.

%==============================================================================
\section{Method: ACR-VLA}
\label{sec:method}
%==============================================================================

We briefly describe the ACR-VLA pipeline to make our ablation experiments interpretable. The pipeline operates on a base VLA that generates $K{=}10$ action samples per step via flow-matching.

\subsection{Conformal Geometry Score (CGS)}

Given $K{=}10$ action samples $\{a_k\}_{k=1}^{K} \in \mathbb{R}^{7}$ (6-DOF end-effector pose + gripper), CGS computes:
\begin{enumerate}[nosep]
\item \textbf{Covariance estimation}: Ledoit-Wolf shrinkage estimator $\hat{\Sigma} = (1-\lambda)S + \lambda \mu I$, where $S$ is the sample covariance, $\mu = \mathrm{tr}(S)/p$, and $\lambda$ is the optimal shrinkage intensity.
\item \textbf{Per-DOF features} (7D): marginal standard deviations from $\mathrm{diag}(\hat{\Sigma})^{1/2}$.
\item \textbf{Per-DOF conformal radii} (7D): $r_j = Q_{1-\alpha}(|a_{k,j} - \bar{a}_j| / \hat{\sigma}_j)$ for each DOF $j$.
\item \textbf{Cross-DOF covariance} (21D): unique off-diagonal entries of the correlation matrix.
\item \textbf{Ellipsoid features} (3D): log-volume, eccentricity, and condition number of the Mahalanobis ellipsoid.
\end{enumerate}
Total: 45D structured conformal features per timestep, concatenated with a 128D observation embedding as input to CRM. In comparison, ReconVLA uses a 1D Euclidean distance score.

\subsection{Adaptive Conformal Inference (ACI) with COC}

ACI adapts the conformal threshold $\alpha_t$ online using Chunk Overlap Consistency (COC) as a self-referential coverage signal:
\begin{equation}
\alpha_{t+1} = \alpha_t + \gamma (\mathrm{COC}_t - (1-\alpha_{\mathrm{target}})),
\end{equation}
where COC measures agreement between overlapping action chunks from consecutive predictions. The target coverage is set to $1-\alpha_{\mathrm{target}} = 0.90$.

\subsection{Conformal Residual Module (CRM)}

CRM is a lightweight MLP (131K parameters) that maps the concatenation of observation embedding (128D) and conformal features (45D) to an action residual $\delta_a \in \mathbb{R}^7$:
\begin{equation}
a_{\mathrm{corrected}} = \bar{a} + \mathrm{CRM}(o, f_{\mathrm{conf}}),
\end{equation}
where $\bar{a}$ is the mean of $K$ samples and $o$ is the observation embedding. CRM is trained on LIBERO demonstrations with synthetic augmentation.

%==============================================================================
\section{Experiments}
\label{sec:experiments}
%==============================================================================

All experiments use SmolVLA as the base VLA, evaluated on LIBERO (nominal, 50 tasks) and LIBERO-Plus (levels L3--L5, 7 perturbation types). We report success rates averaged over 3 seeds unless noted. Statistical significance uses paired $t$-tests across task sets.

%------------------------------------------------------------------------------
\subsection{Full Pipeline Evaluation}
\label{sec:exp-full}
%------------------------------------------------------------------------------

Table~\ref{tab:main} presents the full ACR-VLA pipeline results.

\begin{table}[h]
\centering
\caption{LIBERO / LIBERO-Plus success rates. ACR-VLA provides modest gains on LIBERO-Plus but degrades nominal performance.}
\label{tab:main}
\begin{tabular}{lccccc}
\toprule
Method & Nominal & L3 & L4 & L5 & Plus-Avg \\
\midrule
SmolVLA (base) & 95.0\% & 42.0\% & 28.0\% & 15.0\% & 28.3\% \\
ACR-VLA (full) & 91.0\% & 45.0\% & 31.0\% & 18.0\% & 31.3\% \\
\bottomrule
\end{tabular}
\end{table}

ACR-VLA achieves +3.0 percentage points (pp) on LIBERO-Plus average but incurs a 4.2pp nominal performance drop. The nominal degradation signals that CRM over-corrects in-distribution samples where the base policy is already accurate. CGS covariance diagnostics (Table~\ref{tab:cov}) reveal the root cause of limited gains.

\begin{table}[h]
\centering
\caption{CGS covariance diagnostics: MVP (3D) vs.\ VLA (7D). High shrinkage intensity and low effective rank indicate covariance degeneracy in the VLA setting.}
\label{tab:cov}
\begin{tabular}{lcc}
\toprule
Metric & MVP (3D) & VLA (7D) \\
\midrule
Ledoit-Wolf shrinkage intensity & 0.42 & 0.91 \\
Effective rank of $\hat{\Sigma}$ & 2.7 & 1.3 \\
Off-diagonal magnitude (mean $|r_{ij}|$) & 0.38 & 0.06 \\
Condition number (shrunk $\hat{\Sigma}$) & 4.2 & 1.8 \\
\bottomrule
\end{tabular}
\end{table}

With $K{=}10$ samples in 7D ($p/n = 0.7$), Ledoit-Wolf shrinkage intensity reaches 0.91, effectively collapsing the estimated covariance to a scaled identity matrix. The effective rank drops from 2.7 (MVP) to 1.3 (VLA), and mean off-diagonal correlation drops from 0.38 to 0.06. This means CGS per-DOF features degenerate to near-isotropic, dimension-independent signals, eliminating the multivariate advantage that was observed in 3D settings.

%------------------------------------------------------------------------------
\subsection{Ablation A3: Conformal Features Kill Test}
\label{sec:exp-a3}
%------------------------------------------------------------------------------

The most critical ablation removes all 45D conformal features from CRM, leaving only the 128D observation embedding---functionally equivalent to a generic residual policy (A2C2-style).

\begin{table}[h]
\centering
\caption{A3 ablation: removing conformal features. The 0.8pp gap is statistically insignificant ($p{=}0.38$), triggering the pre-registered kill criterion ($<$2pp).}
\label{tab:a3}
\begin{tabular}{lcccc}
\toprule
Method & Nominal & Plus-Avg & $\Delta$ vs Full & $p$-value \\
\midrule
ACR-VLA (full, 45D conf.) & 91.0\% & 31.3\% & --- & --- \\
A3: no conformal (obs.\ only) & 92.0\% & 30.5\% & $-0.8$ & 0.38 \\
SmolVLA (base) & 95.0\% & 28.3\% & --- & --- \\
\bottomrule
\end{tabular}
\end{table}

\begin{table}[h]
\centering
\caption{A3 per-level breakdown. Conformal features contribute $\leq$1pp at every perturbation level.}
\label{tab:a3-level}
\begin{tabular}{lccc}
\toprule
Level & Full ACR-VLA & A3 (no conf.) & $\Delta$ \\
\midrule
L3 & 45.0\% & 44.5\% & $-0.5$ \\
L4 & 31.0\% & 30.0\% & $-1.0$ \\
L5 & 18.0\% & 17.0\% & $-1.0$ \\
\bottomrule
\end{tabular}
\end{table}

\textbf{Finding 1:} The entire ACR-VLA pipeline gain ($\sim$3pp over base) comes from observation-conditioned residual correction, not conformal features. The A3 variant (no conformal input) still achieves 30.5\% vs.\ the base's 28.3\%, capturing most of the improvement. Conformal features contribute at most 0.8pp---within noise. This directly refutes the claim that per-DOF conformal statistics provide actionable information to the correction network in the $K{=}10$, 7D regime.

%------------------------------------------------------------------------------
\subsection{Ablation A4: Per-DOF vs.\ Scalar Uncertainty}
\label{sec:exp-a4}
%------------------------------------------------------------------------------

\begin{table}[h]
\centering
\caption{A4 ablation: per-DOF (45D) vs.\ scalar (1D, Euclidean distance) vs.\ none (0D). All variants are statistically indistinguishable.}
\label{tab:a4}
\begin{tabular}{lcccc}
\toprule
Feature Type & Dim & Plus-Avg & $\Delta$ vs Full & $p$-value \\
\midrule
Per-DOF (full ACR-VLA) & 45D & 31.3\% & --- & --- \\
Scalar (Euclidean) & 1D & 31.0\% & $-0.3$ & 0.62 \\
No conformal (A3) & 0D & 30.5\% & $-0.8$ & 0.38 \\
\bottomrule
\end{tabular}
\end{table}

\textbf{Finding 2:} Per-DOF granularity provides zero benefit over a single scalar uncertainty score. Combined with A3, this forms a complete negative evidence chain: conformal features contribute negligibly (A3), and within that negligible contribution, per-DOF structure adds nothing beyond scalar (A4). In the MVP 3D setting, per-DOF features showed $\sim$17\% relative improvement over scalar; this advantage vanishes entirely when covariance degeneracy eliminates cross-dimensional structure.

%------------------------------------------------------------------------------
\subsection{Ablation A5: CGS vs.\ Bonferroni}
\label{sec:exp-a5}
%------------------------------------------------------------------------------

\begin{table}[h]
\centering
\caption{A5: joint Mahalanobis (CGS) vs.\ per-DOF Bonferroni correction. No significant difference.}
\label{tab:a5}
\begin{tabular}{lccc}
\toprule
CP Method & Plus-Avg & Ellipsoid Vol.\ (log) & $\Delta$ \\
\midrule
CGS (joint Mahalanobis) & 31.3\% & 2.31 & --- \\
Bonferroni (per-DOF) & 31.1\% & 2.45 & $-0.2$ \\
\bottomrule
\end{tabular}
\end{table}

When shrinkage pushes $\hat{\Sigma}$ toward diagonal form, the joint Mahalanobis ellipsoid and the per-DOF Bonferroni hyperrectangle become geometrically similar (both near-axis-aligned). The 0.2pp gap ($p{=}0.74$) confirms that multivariate CP offers no advantage over marginal corrections in this regime.

%------------------------------------------------------------------------------
\subsection{Ablation A7: Epistemic vs.\ Aleatoric Features}
\label{sec:exp-a7}
%------------------------------------------------------------------------------

\begin{table}[h]
\centering
\caption{A7: CGS epistemic (multi-sample) vs.\ DiU aleatoric (single-pass) features. All variants are indistinguishable, confirming that uncertainty feature type is irrelevant when $K{=}10$ covariance is degenerate.}
\label{tab:a7}
\begin{tabular}{lccc}
\toprule
Feature Source & Dim & Plus-Avg & $\Delta$ vs CGS \\
\midrule
CGS epistemic & 45D & 31.3\% & --- \\
DiU aleatoric & 7D & 31.6\% & $+0.3$ \\
Combined & 52D & 31.4\% & $+0.1$ \\
No features (A3) & 0D & 30.5\% & $-0.8$ \\
\bottomrule
\end{tabular}
\end{table}

DiU aleatoric features perform marginally (but insignificantly) better than CGS epistemic features, likely because they do not depend on multi-sample covariance estimation. The combined 52D variant shows no improvement, indicating both feature types are either redundant with or ignored in favor of observation embeddings.

%------------------------------------------------------------------------------
\subsection{Feature Importance Analysis}
\label{sec:exp-feature}
%------------------------------------------------------------------------------

\begin{table}[h]
\centering
\caption{CRM gradient magnitude by feature group. Observation features dominate; conformal features are $14\times$ smaller with near-zero variance, confirming CRM learns to ignore them.}
\label{tab:grad}
\begin{tabular}{lccc}
\toprule
Feature Group & Mean $|\nabla|$ & Std $|\nabla|$ & Top-5 \\
\midrule
Observation embedding (128D) & 0.041 & 0.018 & 5/5 \\
Per-DOF std (7D) & 0.004 & 0.001 & 0/5 \\
Per-DOF conformal radius (7D) & 0.003 & 0.001 & 0/5 \\
Ellipsoid features (3D) & 0.002 & 0.001 & 0/5 \\
Cross-DOF covariance (28D) & 0.002 & 0.001 & 0/5 \\
\bottomrule
\end{tabular}
\end{table}

\textbf{Finding 3:} Gradient analysis reveals the mechanism behind the ablation results: CRM assigns $14\times$ more gradient magnitude to observation features than conformal features. All top-5 features by importance come from the observation embedding. Crucially, no perturbation-specific pattern exists---CRM does not selectively attend to different DOF features under different perturbation types, directly contradicting the dimension-selective correction hypothesis.

%------------------------------------------------------------------------------
\subsection{COC Signal Validation and ACI Analysis}
\label{sec:exp-coc}
%------------------------------------------------------------------------------

\begin{table}[h]
\centering
\caption{COC signal quality metrics. The correlation with task success is weak ($\rho{=}0.18$), failing the $>$0.3 threshold.}
\label{tab:coc}
\begin{tabular}{lccc}
\toprule
Metric & Value & Threshold & Pass? \\
\midrule
Spearman $\rho$ (COC vs.\ success) & 0.18 & $>0.3$ & Fail \\
$p$-value & 0.04 & $<0.05$ & Pass (weak $\rho$) \\
Ellipsoid narrowing vs.\ static & 6.5\% & $>10\%$ & Fail \\
ACI coverage (target = 0.90) & 0.89 & $\pm 0.02$ & Pass \\
\bottomrule
\end{tabular}
\end{table}

\begin{table}[h]
\centering
\caption{ACI convergence by perturbation level. Convergence rate degrades substantially at higher perturbation levels.}
\label{tab:aci}
\begin{tabular}{lccc}
\toprule
Level & Converged \% & Median Steps & Final Coverage \\
\midrule
L3 & 82\% & 62 & 0.90 \\
L4 & 70\% & 91 & 0.88 \\
L5 & 58\% & 118 & 0.87 \\
\bottomrule
\end{tabular}
\end{table}

\textbf{Finding 4:} COC is a weak proxy for task success. Spearman $\rho = 0.18$ is statistically significant (due to large $N$) but practically negligible. COC measures \emph{prediction consistency}, not \emph{prediction correctness}---a consistently wrong policy produces high COC but low success. ACI achieves near-target coverage (0.89 vs.\ 0.90) via CP's theoretical guarantee, but ellipsoid narrowing is only 6.5\% versus static calibration, indicating ACI's online adaptation provides negligible tightening. Convergence degrades at higher perturbation levels: only 58\% of L5 episodes converge, with the remainder never reaching stable $\alpha$.

%------------------------------------------------------------------------------
\subsection{Ablations A1 and A2: Isolating ACI and CRM}
\label{sec:exp-a1a2}
%------------------------------------------------------------------------------

\begin{table}[h]
\centering
\caption{A1 (no ACI, static calibration) and A2 (no CRM, sample selection only).}
\label{tab:a1a2}
\begin{tabular}{lccc}
\toprule
Method & Nominal & Plus-Avg & $\Delta$ vs Base \\
\midrule
SmolVLA (base) & 95.0\% & 28.3\% & --- \\
A2: CGS+ACI only (no CRM) & 94.0\% & 29.5\% & $+1.2$ \\
A1: Static CGS (no ACI) & 91.5\% & 30.8\% & $+2.5$ \\
Full ACR-VLA (ACI + CRM) & 91.0\% & 31.3\% & $+3.0$ \\
\bottomrule
\end{tabular}
\end{table}

\textbf{Finding 5:} CRM (residual correction) is the only component with meaningful contribution. Without CRM (A2), the pipeline reduces to multivariate CP-based sample selection, gaining only 1.2pp over base---comparable to ReconVLA's scalar approach, confirming that multivariate sample selection offers little over Euclidean distance when covariance is degenerate. ACI contributes only 0.5pp beyond static calibration (A1 vs.\ full), consistent with the weak COC signal analysis.

%------------------------------------------------------------------------------
\subsection{Ablation A6: COC vs.\ Episode-Level Signal for ACI}
\label{sec:exp-a6}
%------------------------------------------------------------------------------

\begin{table}[h]
\centering
\caption{A6: ACI with self-referential COC vs.\ episode-level success signal. The oracle signal provides better ellipsoid narrowing.}
\label{tab:a6}
\begin{tabular}{lccc}
\toprule
Signal Source & Plus-Avg & Ellipsoid Narrowing & External? \\
\midrule
COC (self-referential) & 31.3\% & 6.5\% & No \\
Episode success & 31.8\% & 12.8\% & Yes \\
Static (no adaptation) & 30.8\% & 0\% & No \\
\bottomrule
\end{tabular}
\end{table}

\textbf{Finding 6:} Episode-level success signal outperforms COC for ACI adaptation (12.8\% vs.\ 6.5\% ellipsoid narrowing), despite being coarser-grained and requiring external supervision. This inverts the intended value proposition of self-referential calibration. The practical implication: when external reward signal is available (e.g., in simulation or with success detectors), it should be preferred over consistency-based proxies.

%------------------------------------------------------------------------------
\subsection{SRPO Stacking}
\label{sec:exp-srpo}
%------------------------------------------------------------------------------

\begin{table}[h]
\centering
\caption{Stacking ACR-VLA on SRPO (training-time RL method). Gain is minimal and not significant.}
\label{tab:srpo}
\begin{tabular}{lccc}
\toprule
Method & Plus-Avg & $\Delta$ vs SRPO & $p$-value \\
\midrule
SRPO (baseline) & 82.1\% & --- & --- \\
SRPO + ACR-VLA & 83.2\% & $+1.1$ & 0.29 \\
\bottomrule
\end{tabular}
\end{table}

\textbf{Finding 7:} Inference-time correction provides diminishing returns atop a strong training-time method. SRPO already closes most of the robustness gap (28.3\% $\to$ 82.1\%), leaving primarily structural failures (unseen layout combinations) that action residuals cannot address. Furthermore, SRPO's RL-optimized policy produces more peaked action distributions, reducing flow-matching sample diversity and further degrading CGS quality.

%------------------------------------------------------------------------------
\subsection{Per-Perturbation Breakdown}
\label{sec:exp-perturb}
%------------------------------------------------------------------------------

\begin{table}[h]
\centering
\caption{ACR-VLA gain by perturbation type. Only perturbations coverable by synthetic augmentation (camera, noise) show significant improvement; layout and language exhibit negative transfer.}
\label{tab:perturb}
\begin{tabular}{lcccc}
\toprule
Perturbation & Base & ACR-VLA & $\Delta$ & Significant? \\
\midrule
Camera & 38.0\% & 43.2\% & $+5.2$ & Yes ($p<0.05$) \\
Noise & 35.0\% & 39.8\% & $+4.8$ & Yes ($p<0.05$) \\
Robot-state & 30.0\% & 33.1\% & $+3.1$ & Marginal ($p{=}0.08$) \\
Light & 42.0\% & 44.5\% & $+2.5$ & No ($p{=}0.15$) \\
Background & 40.0\% & 41.8\% & $+1.8$ & No ($p{=}0.22$) \\
Layout & 22.0\% & 21.5\% & $-0.5$ & No \\
Language & 18.0\% & 16.8\% & $-1.2$ & No \\
\bottomrule
\end{tabular}
\end{table}

\textbf{Finding 8:} Only 2 of 7 perturbation dimensions show statistically significant gains, both coverable by standard data augmentation (random crop/resize, Gaussian noise). Layout and language---perturbations requiring semantic understanding beyond pixel-level augmentation---show negative transfer: CRM's learned residuals are counterproductive on out-of-augmentation-distribution inputs. This reveals a fundamental limitation: inference-time correction methods are bounded by the coverage of their training-time augmentation strategy.

%------------------------------------------------------------------------------
\subsection{SnapFlow Integration}
\label{sec:exp-snapflow}
%------------------------------------------------------------------------------

\begin{table}[h]
\centering
\caption{SnapFlow 1-step distillation achieves $7.9\times$ speedup but collapses sample diversity.}
\label{tab:snap}
\begin{tabular}{lccc}
\toprule
Configuration & Inference (ms) & Sample Diversity & Plus-Avg \\
\midrule
Original (50-step) & 142 & 0.023 & 31.3\% \\
SnapFlow (1-step) & 18 & 0.008 & 29.2\% \\
\bottomrule
\end{tabular}
\end{table}

SnapFlow achieves $7.9\times$ inference speedup but preserves only 35\% of sample diversity (std: 0.023 $\to$ 0.008). Since multi-sample CP methods fundamentally depend on sample diversity for covariance estimation, aggressive distillation further degrades an already insufficient statistical foundation. The 2.1pp performance drop, while moderate, compounds with the existing covariance degeneracy problem.

%==============================================================================
\section{Discussion: When Does Multivariate CP Help?}
\label{sec:discussion}
%==============================================================================

Our results identify a clear dimensionality-driven transfer failure. We can characterize the regime boundary:

\paragraph{The $p/n$ ratio threshold.}
In the MVP 3D setting, $K{=}10$ samples give $p/n = 0.3$, where Ledoit-Wolf shrinkage intensity is 0.42 and effective rank is 2.7---sufficient to preserve meaningful covariance structure. In the VLA 7D setting, $p/n = 0.7$ pushes shrinkage to 0.91 and effective rank to 1.3. The transition from ``useful'' to ``degenerate'' covariance occurs roughly when shrinkage intensity exceeds $\sim$0.7 (effective rank $< 2$).

\paragraph{Implications for practitioners.}
\begin{enumerate}[nosep]
\item \textbf{Check covariance quality before building on it.} Ledoit-Wolf shrinkage intensity $> 0.7$ indicates that multivariate CP will not outperform marginal methods. Increase $K$ or reduce action dimensionality.
\item \textbf{Consistency $\neq$ correctness.} Self-referential calibration signals (like COC) measure prediction stability, not accuracy. Under high distribution shift, consistently wrong predictions inflate false confidence.
\item \textbf{Match correction to augmentation coverage.} Inference-time residual correction benefits are confined to the perturbation types seen during training augmentation. On out-of-coverage perturbations, correction can be harmful.
\item \textbf{Diminishing returns atop strong base policies.} When training-time methods have already closed most of the robustness gap, inference-time correction addresses only the structural tail---which typically requires semantic rather than geometric intervention.
\end{enumerate}

\paragraph{Path forward.}
The $K{=}10$ limitation is not fundamental to the approach but to current VLA architectures: flow-matching generation cost scales linearly with $K$. If future architectures reduce per-sample cost (via amortized generation or latent-space sampling), $K \geq 30$ becomes practical, potentially restoring the $p/n$ ratio to a useful regime. Alternatively, dimensionality reduction of the action space (e.g., via learned skill primitives) could reduce $p$ rather than increasing $K$.

%==============================================================================
\section{Conclusion}
\label{sec:conclusion}
%==============================================================================

We conducted a systematic empirical study of multivariate conformal prediction for VLA action correction. Across 11 experiments, we found that the per-DOF advantage observed in low-dimensional settings does not transfer to practical VLA configurations ($K{=}10$, 7D action space). The root cause is covariance degeneracy from insufficient samples relative to dimensionality, quantifiable via Ledoit-Wolf shrinkage intensity.

Our negative results are actionable: they establish that multivariate CP for action correction requires either more samples ($K \geq 30$) or lower-dimensional action representations, and that self-referential calibration signals should not be trusted as proxies for correction quality under significant distribution shift.

The observation-conditioned residual module (CRM without conformal features) does provide modest robustness improvement ($+2.2$pp), but this contribution is functionally equivalent to existing residual policy methods and does not constitute a novel contribution.

We believe systematic negative results, accompanied by precise failure analysis, are valuable for guiding the community away from unproductive directions and toward the conditions where structured uncertainty quantification can genuinely improve robotic manipulation.

%==============================================================================
% References
%==============================================================================
\bibliographystyle{plainnat}

\begin{thebibliography}{10}

\bibitem[Vovk et al., 2005]{vovk2005}
V.~Vovk, A.~Gammerman, and G.~Shafer.
\newblock {\em Algorithmic Learning in a Random World}.
\newblock Springer, 2005.

\bibitem[Ledoit and Wolf, 2004]{ledoit2004}
O.~Ledoit and M.~Wolf.
\newblock A well-conditioned estimator for large-dimensional covariance matrices.
\newblock {\em Journal of Multivariate Analysis}, 88(2):365--411, 2004.

\bibitem[Kim et al., 2024]{kim2024}
M.~Kim et al.
\newblock LIBERO-Plus: Robustness benchmarks for vision-language-action models.
\newblock 2024.

\bibitem[Luo et al., 2024]{luo2024}
J.~Luo et al.
\newblock ReconVLA: Uncertainty-aware sample selection for robust VLA deployment.
\newblock 2024.

\bibitem[Ren et al., 2024]{ren2024}
A.~Ren et al.
\newblock SRPO: Self-refined policy optimization for robust VLA.
\newblock 2024.

\bibitem[Gibbs and Cand\`es, 2021]{gibbs2021}
I.~Gibbs and E.~Cand\`es.
\newblock Adaptive conformal inference under distribution shift.
\newblock {\em NeurIPS}, 2021.

\bibitem[Black et al., 2024]{black2024}
K.~Black et al.
\newblock $\pi_0$: A vision-language-action flow model for general robot control.
\newblock 2024.

\bibitem[Liang et al., 2024]{liang2024}
Z.~Liang et al.
\newblock A2C2: Action-aware conformal correction for robotic manipulation.
\newblock 2024.

\bibitem[Dixit et al., 2024]{dixit2024}
A.~Dixit et al.
\newblock SAFE: Slow and fast parameter-efficient tuning for continual learning.
\newblock 2024.

\bibitem[Wu et al., 2024]{wu2024}
Y.~Wu et al.
\newblock DiG-Flow: Diversity-guided flow matching for robust VLA.
\newblock 2024.

\end{thebibliography}

\end{document}
